% !PRINTED 428
\weberpage{428}
Die Differenz
\[
x(x-1)(x-2)\dots(x-n+1) - x^n + x
\]
ist eine Function, h\"ochstens vom $n-1^{\text{ten}}$ Grade. Sie ist aber
f\"ur $n$ Werthe von $x$, n\"amlich f\"ur $x = 0, 1, 2, \dots n-1$ congruent mit Null, und daher haben wir identisch
\[
x(x-1)(x-2)\dots(x-n+1) \equiv x^n - x \quad (\mathrm{mod}\ n).
\]
Daraus ergiebt sich der \gesperrt{Wilson'sche Lehrsatz}:
\[
1\,.\,2\,.\,3\dots(n-1) \equiv -1 \quad (\mathrm{mod}\ n),
\]
wenn man die Co\"efficienten der ersten Potenz von $x$ beiderseits
vergleicht.

Wir beschr\"anken uns jetzt auf die zweite der Congruenzen (5)
\weq{8}{x^{n-1} \equiv 1 \quad (\mathrm{mod}\ n)}
und beweisen zun\"achst den Satz:

\gesperrt{3. Ist $a$ ein Theiler von $n-1$, so hat die Congruenz}
\weq{9}{x^a \equiv 1 \quad (\mathrm{mod}\ n)}
\gesperrt{genau $a$ verschiedene Wurzeln.}

Denn es ist, wenn $n-1 = ab$ ist,
\[
x^{n-1} - 1 = (x^a - 1)(x^{a(b-1)} + x^{a(b-2)} + \cdots + x^a + 1),
\]
und da jede Wurzel der linken Seite Wurzel entweder des einen
oder des anderen Factors auf der rechten Seite sein muss, so
folgt, dass, wenn $x^a - 1 \equiv 0$ weniger als $a$ Wurzeln h\"atte, der
zweite Factor mehr Wurzeln haben m\"usste, als sein Grad angiebt,
entgegen dem Satz 1.

Ist $\alpha$ eine durch $n$ nicht theilbare Zahl, und $a$ die kleinste
positive Zahl, f\"ur die $\alpha^a \equiv 1 \pmod n$ wird, ist ferner $h$ irgend
eine positive ganze Zahl, f\"ur die $\alpha^h \equiv 1 \pmod n$ ist, so ist $a$
nothwendig ein Theiler von $h$, insbesondere also $a$ ein Theiler
von $n - 1$.

Denn setzen wir $h = qa + a'$, worin $q$ eine ganze Zahl und
$a' < a$ ist, so ist auch $\alpha^{a'} \equiv 1$; also muss $a' = 0$ sein, weil $a$ die
kleinste positive Zahl sein sollte, wof\"ur diese Congruenz befriedigt ist.

\gesperrt{Eine Wurzel $\alpha$ der Congruenz (9) von der
Eigenschaft, dass keine niedrigere als die $a^{\text{te}}$
Potenz der Einheit congruent wird, heisst eine
primitive Wurzel dieser Congruenz, oder primitive $a^{\text{te}}$ Congruenzwurzel der Primzahl $n$, und eine
primitive Wurzel $g$ der Congruenz (8) nennen wir
auch kurz eine primitive Wurzel der Primzahl $n$.}
